% 2026-10-01: authorized speech reorganization for 26 main slides and five hidden backups.
% Spoken bullets match the current PowerPoint notes. Slide content and media are preserved.
\documentclass[12pt,a4paper]{article}
\usepackage[T1]{fontenc}
\usepackage[utf8]{inputenc}
\usepackage[english]{babel}
\usepackage[a4paper,margin=22mm]{geometry}
\usepackage{amsmath,amssymb}
\IfFileExists{newtxtext.sty}{\usepackage{newtxtext,newtxmath}}{\usepackage{mathptmx}}
\usepackage{microtype}
\usepackage{xcolor}
\usepackage{setspace}
\usepackage{enumitem}
\usepackage[hidelinks,unicode]{hyperref}
\definecolor{navy}{HTML}{183B63}
\hypersetup{pdftitle={Thesis Defense -- Speaking Script},pdfauthor={Heykel Khadhraoui}}
\setstretch{1.12}
\setlength{\parindent}{0pt}
\setlength{\parskip}{0pt}
\setlength{\emergencystretch}{2em}
\widowpenalty=10000
\clubpenalty=10000
\pagestyle{empty}
\makeatletter
\renewcommand\section{\@startsection{section}{1}{0pt}{14pt}{6pt}{\normalfont\Large\bfseries\color{navy}}}
\makeatother
\setlist[itemize]{leftmargin=1.35em,label=\textcolor{navy}{\textbullet},itemsep=2pt,parsep=0pt,topsep=3pt,partopsep=0pt}
\begin{document}

\section*{Master Thesis Presentation}
\begin{itemize}
\item Hello everyone. Today, I will be presenting what I have been working on over the past few months.
\item I developed a surface grinding controller for a seven-degree-of-freedom Franka robot. The control method I used is Cartesian impedance control.
\end{itemize}

\section*{Motivation}
\begin{itemize}
\item The main problem in surface grinding is that we cannot ensure that the configured and physical surfaces are perfectly aligned. Calibration leaves some uncertainty in the surface orientation.
\item Joint friction can also prevent the tool from reaching the desired orientation. The tool may therefore first contact the surface at an edge.
\item The idea is to use impedance control to make rotation compliant. Pressing the tilted tool creates a contact moment, which can rotate the tool towards alignment.
\end{itemize}

\section*{Contact demonstration}
\begin{itemize}
\item The video shows pickup, approach, contact, and the hold before grinding.
\item During the hold, I change the surface orientation. Normal pressing continues as contact moments rotate the tool.
\item The desired orientation stays fixed. This is a qualitative demonstration of compliant contact.
\end{itemize}

\section*{Overview}
\begin{itemize}
\item After this introduction, I explain the Cartesian controller, followed by the contact experiments and results.
\item The separate null-space controller, experiment and results follow together. I finish with the conclusion and future work.
\end{itemize}

\section*{Cartesian pose}
\begin{itemize}
\item To begin, let's first introduce Cartesian pose. Cartesian pose describes the position and orientation of the gripper, also called the end effector, or EE.
\item It has three positional degrees of freedom and three rotational degrees of freedom. The pose is determined through forward kinematics from the positions of all seven joints.
\end{itemize}

\section*{Surface frame}
\begin{itemize}
\item To choose compliance relative to the surface, we use a surface frame. The normal points out of the surface, while the two perpendicular tangents lie in the surface.
\item This lets us choose different stiffnesses for pressing, tangential motion and rotation. I will use these directions when explaining the contact settings.
\end{itemize}

\section*{Cartesian impedance controller}
\begin{itemize}
\item Impedance control uses a spring-damper model that produces a force to correct position and a moment to correct rotation. Together, the force and moment form what is called a wrench.
\item K and D are the stiffness and damping matrices. In the surface frame they are diagonal, so translation and rotation are decoupled in this form.
\end{itemize}

\section*{Real-time control}
\begin{itemize}
\item The controller compares the reference state with the measured state. Here, state includes both pose and velocity, so stiffness and damping act on their respective errors.
\item The wrench is then mapped to joint torques through the transpose of the Jacobian. The Jacobian itself maps the seven joint velocities to the linear and angular velocity of the end effector.
\[
\tau=J^\top(q)F
\]
\[
\begin{bmatrix}\dot p_{\mathrm{EE}}\\\omega_{\mathrm{EE}}\end{bmatrix}=J(q)\dot q
\]
\item The complete command also includes model compensation and any enabled null-space torques. The controller operates in real time, with an update every millisecond.
\end{itemize}

\section*{Centre of compliance (CoC)}
\begin{itemize}
\item An interesting idea presented in this work is to couple force and moment so that a pressing force also produces a rotational moment. We do this by choosing a virtual reference point called the centre of compliance, or CoC.
\item We can illustrate this behaviour here. In the left panel, the displacement from the tool centre point gives an additional moment that supports alignment.
\item In the right panel, the displacement is on the other side and the additional moment opposes alignment. This motivates the reference-point transformation on the next slide.
\end{itemize}

\section*{Translation--rotation coupling}
\begin{itemize}
\item The tool centre point, or TCP, is the controlled reference point on the end effector. We shift the reference point of the impedance model to the virtual CoC, with the displacement defined from the TCP to that point.
\item The commanded force then acts as if applied at the virtual point. This changes the controller's force and moment response without moving the physical contact point.
\item For the local, small-rotation model, we transform the stiffness and damping matrices using the point-shift adjoint. The resulting matrices contain off-diagonal terms, which couple translation and rotation.
\end{itemize}

\section*{Contact experiment}
\begin{itemize}
\item Our contact experiment proceeds as follows. First, we orient the tool towards the configured surface and slowly approach it along its normal.
\item During contact establishment, compliant rotation allows the contact moment to turn the tool towards alignment. Grinding then begins along the tangential direction \(t_2\).
\item For translation, we choose a compliant normal stiffness to allow pressing against the surface. The tangential directions remain stiff to limit unwanted motion during alignment.
\item For rotation, the normal direction is stiff, while rotations about the tangential directions are compliant. This allows the tool to tilt in response to contact moments.
\end{itemize}

\section*{Plausibility experiment}
\begin{itemize}
\item Before looking at the contact results, we first check whether the controller produces the expected force and moment. This video introduces the plausibility experiment.
\item The next two plots compare the commanded and model-estimated responses with the spring-law predictions for manual translation and rotation.
\end{itemize}

\section*{Normal-force plausibility assessment}
\begin{itemize}
\item First, we manually push the end effector about 20 mm along the surface normal, while keeping the desired pose fixed and the CoC at the TCP. With a normal stiffness of 1000 newtons per metre, the measured position-error change predicts a force of about minus 19.65 newtons.
\[
F_{n,\mathrm{qs}}=(1000\,\mathrm{N/m})(-0.019650\,\mathrm{m})=-19.650\,\mathrm{N}
\]
\item In the shaded stationary interval, the mean commanded force is minus 19.64 newtons, very close to the quasi-static prediction. The model-estimated force is minus 22.29 newtons, about 13.5 percent different from the command.
\item The commanded response therefore agrees with the spring law in this test. The model estimate provides a separate plausibility check.
\end{itemize}

\section*{Moment plausibility assessment}
\begin{itemize}
\item We do the same for the moment by manually rotating the tool about the first surface tangent while retaining the normal push. The measured rotational-error change is about 6.57 degrees relative to that loaded baseline.
\item With a rotational stiffness of 15 newton metres per radian, the prediction is about 1.72 newton metres. The stationary commanded moment matches this value, while the model estimate is about 1.85 newton metres, a difference of roughly 8 percent.
\[
M_{t_1,\mathrm{qs}}=(15\,\mathrm{N\,m/rad})\left(6.565\,\frac{\pi}{180}\,\mathrm{rad}\right)\approx1.719\,\mathrm{N\,m}
\]
\item Together, these two checks support the expected force and moment response for the tested errors.
\end{itemize}

\section*{Commanded and estimated wrench}
\begin{itemize}
\item Next, we compare the commanded and model-estimated wrench during a separate contact trial, with the CoC at the TCP. The upper plot shows normal force and the lower plot shows moment about the first surface tangent.
\item The force curves follow each other closely as pressing increases. During alignment, the estimated moment is much larger than the commanded rotational moment.
\item Later, the force and moment settle and the commanded and estimated values become close. The moment difference during alignment is consistent with an external contact moment driving the compliant rotation.
\end{itemize}

\section*{Effect of rotational stiffness}
\begin{itemize}
\item For the contact results, angular error describes the tool-normal tilt relative to the calibrated surface, about the first surface tangent. It is calculated from measured end-effector orientation, and a smaller magnitude means less error about this tangent.
\item This plot shows that higher rotational stiffness produces a larger remaining angular error. Increasing the stiffness from 5 to 50 newton metres per radian increases the error from about 1.75 to 6.58 degrees.
\item The rotational spring therefore resists the contact-driven alignment more strongly. These points are three-trial means, with error bars showing one sample standard deviation.
\end{itemize}

\section*{Effect of CoC position}
\begin{itemize}
\item Next, this plot shows how the centre of compliance affects the angular error. We moved it from minus 100 to plus 100 mm along the second surface tangent for positive and negative entry tilts.
\item For the positive entry tilt, shown in black, positive displacement supports alignment and reduces the error. Large negative displacement opposes it and leaves a much larger error.
\item For the negative entry tilt, shown in blue, the smallest tested error magnitude occurs near plus 10 mm. At larger positive displacements the error becomes negative and its magnitude grows.
\item The important point is that the effect depends on the entry direction and the chosen displacement. A larger displacement does not automatically give a smaller angular error.
\end{itemize}

\section*{Angular error and interaction wrench}
\begin{itemize}
\item Next, we compare representative time histories at CoC displacements of minus 40 mm, zero and plus 40 mm. The three panels show angular error, model-estimated normal force and model-estimated TCP moment.
\item For this positive entry tilt, the supporting plus-40-mm setting reaches its final-error band at about 2.5 seconds, compared with 3.6 seconds at the TCP. The band means staying within 0.1 degree of each trace's own final value.
\item The opposing minus-40-mm setting leaves the angular error close to its entry value. Its estimated moment continues to build while little alignment occurs.
\item The force traces are broadly similar and settle near minus 80 newtons, while the moment traces differ clearly. The supporting displacement changes the moment response and speeds up the angular response in this comparison.
\end{itemize}

\section*{Null-space controller}
\begin{itemize}
\item Another study in this work analysed null-space behaviour. Our Cartesian pose has six degrees of freedom, while the robot has seven joints, leaving one redundant direction when the Jacobian has full rank. Several joints can move together in this direction without changing the instantaneous Cartesian motion.
\item We introduced a damping torque to oppose this motion.
\item The conditioning torque seeks a configuration with a larger minimum singular value of the Jacobian. At a singularity, the Jacobian loses rank and the end effector loses an instantaneous motion direction. Near a singularity, some motions require very high joint velocities.
\end{itemize}

\section*{Disturbance - Without null-space control to damping}
\begin{itemize}
\item This first video shows the response to a torque disturbance, first without null-space control and then with damping. It gives a qualitative view of how the joints respond.
\end{itemize}

\section*{Disturbance - Conditioning only}
\begin{itemize}
\item This is the continuation of the same demonstration, now with conditioning only. The following experiment gives the quantitative comparison between the null-space settings.
\end{itemize}

\section*{Null-space experiment}
\begin{itemize}
\item To study this case, we held the captured Cartesian pose as the reference and added a commanded torque disturbance equivalent to a virtual force of 20 newtons on link three. Through the Jacobian mapping, this disturbance acted mainly on joint one.
\item With the disturbance applied, we compare four conditions with three trials each: no null-space control torque, conditioning alone, damping alone, and both torques together. The enabled settings are 2 newton metres for conditioning and 2 newton metre seconds per radian for damping.
\item The following plots show the same four-second interval after disturbance onset. The curves are three-trial means, and the shading shows one sample standard deviation.
\end{itemize}

\section*{Cumulative joint motion}
\begin{itemize}
\item Cumulative joint motion accumulates the magnitude of projected motion using all seven joint velocities. It includes movement in both directions, so reversals do not cancel each other.
\item With damping alone, cumulative motion is about 25 percent lower than without null-space torque. Conditioning reduces it further in this comparison.
\item Adding damping to conditioning reduces cumulative motion by about another 51 percent compared with conditioning alone. The combined setting gives the smallest cumulative motion among these four conditions.
\end{itemize}

\section*{Jacobian conditioning}
\begin{itemize}
\item Next, we look at the minimum singular value of the Jacobian. Without null-space torque, and with damping alone, this indicator decreases during the disturbance.
\item Conditioning alone and combined control keep it nearly constant. The small difference between these two curves does not establish a meaningful conditioning advantage.
\end{itemize}

\section*{Joint motion over time}
\begin{itemize}
\item Finally, we look at the measured angle change of joint one, which receives most of the disturbance torque. Damping reduces this motion, and the conditioning and combined settings keep it much smaller.
\item Both panels show the full four seconds, with a larger vertical scale on the right to reveal the small motions. Conditioning alone produces repeated reversals, and adding damping reduces total motion while reversals remain.
\end{itemize}

\section*{Conclusion}
\begin{itemize}
\item To conclude, compliant rotation allowed contact moments to reduce angular error. Higher rotational stiffness left a larger error, and the centre of compliance could either support or oppose the contact moment.
\item The null-space controller reduced redundant joint motion under disturbance. Combining conditioning with damping reduced cumulative motion further than conditioning alone in the main comparison.
\item Future work could use a more rigid tool mount for better measurements. The CoC could also be adapted to support alignment and then returned to the TCP.
\item Thank you for your attention.
\end{itemize}

\section*{Opposing moment}
\begin{itemize}
\item This video illustrates the opposing-moment case. An additional moment that opposes the contact-driven rotation can hinder alignment.
\item It complements the supporting and opposing configurations in the earlier illustration. The contact plots provide the quantitative comparison.
\end{itemize}

\section*{Null-space torques and projector}
\begin{itemize}
\item The kinematic projector selects null-space torque directions. It uses the Jacobian pseudoinverse J-plus and the seven-by-seven identity matrix.
\[
N_\tau=I_7-J^\top J^{+\top}
\]
\item Damping and conditioning add to give the total null-space torque.
\[
\tau_{\mathrm{null}}=\tau_d+\tau_\sigma
\]
\item Damping opposes projected joint velocity and is zero at rest. The damping coefficient sets its strength.
\[
\tau_d=-d_{\mathrm{null}}N_\tau\dot q
\]
\item Conditioning probes both directions along the unit null-space vector and chooses the direction with the larger minimum singular value. The selector is plus one or minus one, and zero within the deadband.
\[
\tau_\sigma=k_\sigma N_\tau(s_\sigma v_7)
\]
\item The conditioning coefficient sets the torque magnitude. Unlike damping, this torque can be nonzero at rest.
\end{itemize}

\section*{Angular quantities}
\begin{itemize}
\item The measured angular offset is the tool-normal tilt at contact entry relative to the calibrated surface. The angular error describes that tilt during contact, resolved about the first surface tangent.
\item The red arc shows the entry offset and the blue arc shows the remaining error at contact end. Both are calculated from measured end-effector orientation and the calibrated tool normal.
\item This is different from the controller's orientation error relative to its held reference. A zero first-tangent component alone does not establish perfect physical alignment, because other components and calibration or mount uncertainty remain.
\end{itemize}

\section*{Baseline angular error}
\begin{itemize}
\item This plot shows the baseline with the CoC at the TCP. The measured entry offsets are about 0.69, plus 9.31 and minus 9.41 degrees.
\item The corresponding remaining angular-error components are about 1.65, 1.75 and 1.41 degrees. The two larger entry tilts are reduced substantially, while the nominal-zero case retains a small error.
\item The response is not perfectly symmetric, and friction, calibration and tool-mount effects may contribute. This plot alone does not separate those contributions.
\end{itemize}

\section*{Sources of angular offset}
\begin{itemize}
\item This sketch separates two sources of angular offset at surface entry. The physical surface can differ from the configured surface, and the actual tool orientation can differ from the desired orientation.
\item The dashed black line shows the desired orientation, while the green line shows the tool face. Together, these differences explain why the tool can meet the physical surface at a tilt.
\end{itemize}

\end{document}
